% Chapter 4 - System Overview
\section{Automated OSM-to-CARLA Map Generation Pipeline}
\label{sec:system_overview}


This chapter describes the automated pipeline that converts OpenStreetMap (OSM) data into CARLA-compatible OpenDRIVE (XODR) maps and spatial tiles. CARLA is a high-fidelity simulator for autonomous driving research, but its OpenDRIVE ingestion is sensitive to structural inconsistencies; therefore, the generation process must target both OpenDRIVE correctness and CARLA-specific failure modes \cite{dosovitskiy2017carla,asam_opendrive_spec}. The proposed pipeline is implemented as an auditable sequence of deterministic stages with explicit contracts and persistent artifacts.

\subsection{Overview and design goals}
The thesis requires multiple pipeline runs and artifact variants from OSM data \cite{openstreetmap} to evaluate the structural domain gap between automatically generated maps and a manually authored reference map of Ingolstadt. The map generation process must satisfy three constraints:

\begin{itemize}
    \item \textbf{Geometric validity} with respect to OpenDRIVE (continuous and well-formed planView and elevation profiles).
    \item \textbf{Semantic completeness} required for CARLA driving and sensor simulation (lanes, lane connectivity, and traffic-relevant semantics).
    \item \textbf{Reproducibility} to enable controlled experiments and fair comparisons across runs.
\end{itemize}

To meet these constraints, the pipeline writes a sequence of intermediate OpenDRIVE snapshots and diagnostic artifacts, enabling ablation studies and failure localization when CARLA import issues occur. Figure~\ref{fig:pipeline_stages} shows the high-level stage sequence from OSM extraction to domain-gap metrics.

\begin{figure}[htbp]
  \centering
  \includegraphics[width=0.96\textwidth]{images/fig_pipeline_stages}
  \caption[High-level OSM-to-CARLA pipeline stage sequence.]{High-level stage sequence of the governed OSM-to-CARLA map generation pipeline. Each stage produces a versioned intermediate artifact; the final output feeds directly into the domain-gap metric computation. Source: author's own diagram.}
  \label{fig:pipeline_stages}
\end{figure}

\subsection{Input data and spatial scope}
The spatial scope is defined by a fixed GPS bounding box in Ingolstadt. The bounding box is used to extract the OSM region deterministically and to align downstream elevation data. The manually authored Ingolstadt OpenDRIVE map is treated as an immutable reference baseline; the pipeline generates one or more automatic variants from OSM using identical spatial bounds to isolate map-construction effects from geographic changes. Because OSM semantics are incomplete and heterogeneous, the pipeline uses conservative, deterministic rules for missing lane or junction annotations and logs all inferred decisions as artifacts.



\subsection{Coordinate reference handling and canonical evaluation frame}
\label{sec:crs_handling}
A recurring source of invalid structural gap measurements is inconsistent georeferencing across OpenDRIVE sources.
The manual Ingolstadt reference maps (Grid0821/0828) are provided in a Transverse Mercator projection consistent with the UTM Zone~32N family (metric coordinates in meters), while automatically generated OSM-derived maps may use a local Transverse Mercator projection anchored to the project bounding box and additionally encode a global translation via the \texttt{<offset>} element.
Directly comparing raw \texttt{(x,y)} coordinates across these frames can therefore produce kilometer-scale errors that are unrelated to actual map differences.

To ensure meaningful structural metrics, this thesis adopts a \emph{canonical evaluation frame}:
all geometric operations (alignment, distance metrics, and tile bounding boxes) are performed after transforming map geometry into EPSG:32632 (WGS84 / UTM Zone~32N), with \texttt{<offset>} applied before projection.
The transformation is implemented via deterministic header parsing (including CDATA and whitespace normalization) and is audited via a georeference provenance record written per run.
After canonicalization, alignment is restricted to rigid SE(2) transforms (rotation + translation, scale fixed to 1) and guarded by a sanity threshold to prevent reporting misleading metrics when overlap is insufficient.

Figure~\ref{fig:crs_normalization} summarizes the normalization and alignment workflow.

\begin{figure}[H]
    \centering
    \includegraphics[width=\textwidth]{images/fig_crs_normalization.png}
    \caption{Coordinate normalization used for valid comparison between manual and automatically generated OpenDRIVE maps. Both sources are parsed (geoReference and optional offset), projected to a canonical CRS (EPSG:32632), and then aligned using a rigid SE(2) transform with a sanity gate.}
    \label{fig:crs_normalization}
\end{figure}

\subsection{Stage-by-stage map construction}
The pipeline is organized into sequential stages. Each stage writes a versioned OpenDRIVE output (XODR) and a structured JSON report so that the full transformation chain remains traceable from input to final OpenDRIVE output.

\subsubsection{Sanitation (Stage 1)}
The sanitation stage normalizes OpenDRIVE metadata (e.g., \texttt{geoReference}) and repairs common formatting issues that can prevent downstream tools from parsing the map. It also generates a cropped ``QA slice'' for rapid inspection. The stage output serves as a stable baseline for subsequent topology and geometry processing.

\subsubsection{Topology lint and optional SUMO repair (Stage 2)}
Automatic conversion from OSM frequently produces junction artifacts, missing connections, and inconsistent road-link structures. The pipeline therefore performs a topology lint pass and can optionally execute a SUMO-based repair workflow to stabilize the road graph \cite{sumo}. A validation report is stored as machine-readable JSON to support reproducible error analysis and comparisons between pipeline versions.

\subsubsection{Topology repair (Stage 3)}
Stage~3 applies deterministic cleanup of road predecessor/successor relationships and junction references. The goal is to ensure that the connectivity graph is well-formed before higher-level semantic enrichment is applied.
The current governed pipeline also canonicalizes junction connectors before geometry freeze: true duplicate connector families are collapsed by \texttt{incomingRoad}/\texttt{exitRoad}, and irregular same-start connector bundles may have their \texttt{planView} geometry normalized without deleting distinct turn semantics. The pass writes a Stage~3 canonicalization report and enforces invariants so that no incoming road loses all valid junction connectivity.

\subsubsection{Enrichment and object injection (Stage 4)}
To support perception studies and qualitative inspection, the pipeline optionally inserts semantic objects and environmental structure. Depending on data availability, this includes buildings derived from OSM footprints and, when available, preprocessed building footprints provided as GeoJSON \cite{rfc7946}, as well as additional realism objects. Enrichment is treated as an additive overlay: it is not allowed to mutate plan-view geometry or lane topology and is later verified by integrity gates.

A photometric augmentation module (\texttt{augmentation/RealismAugmentor}) is also provided for image-level realism during sensor data post-processing. It applies Gaussian noise ($\sigma \in [5, 25]$), horizontal and vertical motion blur (kernel sizes 3, 5, 7), brightness and contrast scaling (contrast $\alpha \in [0.7, 1.4]$, brightness $\beta \in [-40, +40]$), and per-channel colour shift ($\pm 20$). The module is designed to run without a mandatory OpenCV dependency (NumPy fallback) and is safe for HPC multi-processing (no global RNG seeding). It has not been applied to governed thesis experiments; its evaluation as a perceptual domain-gap reduction tool is deferred to future work (see §\ref{sec:future_work}).

An auxiliary 3D enrichment path based on \textsc{OSM2World}~\cite{osm2world} was integrated for mesh-level
scene generation. \textsc{OSM2World} converts the OSM source to an \textsc{OBJ} scene file~\cite{osm2world};
a stack-check conversion for the Ingolstadt area (using a validated OSM extract, SHA-256 confirmed)
produced an approximately 398\,MB \texttt{scene.obj} artifact in approximately 252\,seconds.
A subsequent \textsc{Blender}~\cite{blender}-driven OBJ$\to$\textsc{FBX} stage
converts the mesh to the format required for Unreal Engine asset cooking, enabling optional
import of realistic environmental geometry (buildings, vegetation) into a packaged CARLA world.
Both stages use deterministic caching keyed on SHA-256 hashes of the input OSM file and
configuration, writing a \texttt{osm2world\_status.json} or \texttt{blender\_status.json}
provenance record. This auxiliary path is disabled by default (\texttt{ENABLE\_OSM2WORLD=0})
and does not affect OpenDRIVE generation. Importing OSM2World-derived FBX assets into
a distributable CARLA world additionally requires the full Unreal Engine source build of CARLA,
which is outside the scope of this thesis; only the XODR-based standalone API
\texttt{generate\_opendrive\_world} was used for simulator-dependent evaluation here.
The structural domain-gap analysis in this thesis is therefore performed on the final
OpenDRIVE road-network artifact, and OSM2World-derived mesh enrichments lie outside
the primary comparison scope.

\subsubsection{Elevation application from DEM (Stage 5)}
Because the pipeline supports DEM-backed 3D-capable artifacts, elevation is implemented as a first-class pipeline capability rather than a cosmetic detail. A digital elevation model (DEM) aligned to the GPS bounds can be sampled along roads and encoded in OpenDRIVE elevation profiles, and a smoothing operation is applied to avoid unrealistic high-frequency elevation changes that can break collision meshes or produce implausible slopes.
The elevation source in this thesis is a Copernicus DEM product retrieved via OpenTopography, which provides reproducible access to global DEM datasets and metadata for scientific workflows \cite{opentopography,copernicusdem}.
For the authoritative structural analysis discussed in Chapter~7, however, the run\_11 comparison artifact remained flat at $z=0$ even though later governed artifacts proved that DEM-backed generation is possible. The reported primary structural gap metrics are therefore planar (2D) only, and Section~\ref{sec:results}, especially the paragraph ``Elevation status and CARLA visual QA boundary,'' should be read as the authoritative interpretation for the structural analysis artifact.

\subsubsection{PlanView and continuity smoothing (Stage 6)}
OSM-derived geometry often contains fragmented plan-view elements and heading discontinuities. Stage~6 merges micro-segments, smooths plan-view geometry, clamps pathological curvature, and performs continuity repair to reduce geometric noise while preserving the global road layout. This stage typically reduces the number of geometric primitives and improves simulator stability by preventing discontinuities that trigger CARLA import failures.

\subsubsection{Lane and cross-section generation (Stage 7)}
For CARLA driving simulation and dataset generation, lane semantics must be present and consistent. This stage generates or repairs lane structures, applies lane-offset smoothing, and enforces cross-section continuity constraints to prevent negative widths and discontinuities at laneSection boundaries. A lane overlay visualization and cross-section diagnostics are exported for manual plausibility checks.

\subsubsection{Lane links, markings, and integrity gates (Stage 8)}
The final stage creates or repairs lane-to-lane connectivity, generates lane markings, and enforces CARLA-specific successor invariants at the laneSection level. Additional integrity gates (e.g., collision-mesh plausibility checks) detect major polygon issues that are strongly correlated with CARLA import failures. The output of this stage is the \emph{contract-checked} OpenDRIVE file used in subsequent experiments.
After late junction-link mutation, the selected final XODR is rechecked by lane-connectivity, junction-integrity, and connector-crowding gates. Conservative connector pruning and the \texttt{junction\_connector\_risk.json} audit are part of this finalization path; if residual connector crowding exceeds configured thresholds, promotion fails rather than silently accepting the artifact.

\subsection{Quality gates and integrity contracts}
\begin{table}[htbp]
\centering
\small
\begin{tabular}{p{3.0cm}p{3.4cm}p{3.0cm}p{4.2cm}}
\hline
Gate Name & Condition & Action on Failure & Protects Against \\
\hline
Min XODR file size & $\geq 10{,}000$ bytes & Hard stop & Stub fixture inputs (run\_02: 796 bytes) \\
Tile IoU gate & IoU $\geq 0.5$ & Skip per-tile metrics & Invalid local comparisons \\
CRS sanity gate & Overlap after alignment $>$ threshold & Reject, report deferred & CRS mis-projection artefacts \\
Empty correspondence gate & \texttt{correspondence.csv} $\geq$\,1 data row & Hard stop & Silent zero-row claims \\
Claim boundary annotation & \texttt{\_claim\_boundary} in \mbox{\texttt{full\_report.json}} & Required for admissibility & Unannotated result propagation \\
Connector-crowding gate & Repeated-anchor and overlap counts within thresholds & Hard stop & Exploded junction patches / connector crowding \\
\hline
\end{tabular}
\caption{Thesis integrity gates applied to prevent false-positive structural claims.}
\label{tab:quality_gates_contract}
\end{table}

Table~\ref{tab:quality_gates_contract} summarizes the thesis-specific gate conditions that prevent broken correspondence, CRS mismatch, or stub artifacts from being misreported as valid structural evidence.

The run\_02 stub defect (796-byte XODR) was retroactively blocked by Gate~1 after discovery (Table~\ref{tab:quality_gates_contract}); run\_03 was the first re-run with real inputs (15.3 MB auto, 65.9 MB manual XODR), but the final authoritative structural thesis metrics are reported from the governed run\_11 raw basis (Section~\ref{sec:results}).

All gate decisions, override flags, and bypass history are recorded in a machine-readable
\texttt{kill\_switch\_provenance} block within \texttt{full\_report.json}.
Both the identity-alignment guard and the empty-correspondence guard are enforced by default;
they can only be bypassed by explicitly setting \texttt{UP\_ALLOW\_}\allowbreak\texttt{IDENTITY\_ALIGNMENT=1}
or \texttt{UP\_ALLOW\_EMPTY\_}\allowbreak\texttt{CORRESPONDENCE=1} respectively, and any such override is
permanently logged in the provenance block.
Domain-gap metrics are normalized to $[0,1]$ using explicit reference scales drawn from
\texttt{SETTINGS.}\allowbreak\texttt{DOMAIN\_GAP\_}\allowbreak\texttt{NORMALIZATION}: geometry RMSE reference $= 1.0\,\text{m}$,
Hausdorff reference $= 5.0\,\text{m}$, curvature KL reference $= 0.5$,
intersection delta reference $= 1.0$, semantic delta reference $= 1.0$.
The effective normalization applied to each metric is reported verbatim in the
\texttt{normalization\_contract} block of \texttt{full\_report.json}.

\subsection{Outputs and reproducibility artifacts}
Each pipeline run produces (i) a sequence of stage XODR outputs, (ii) per-stage preview images for qualitative verification, and (iii) diagnostic JSON files that quantify structural properties (e.g., road counts, junction counts, curvature statistics) and validation outcomes. A settings snapshot is written to record all run-defining parameters; combined with cryptographic hashes of key artifacts, this snapshot enables reproducible re-runs and supports determinism auditing.




\subsection{Design and reproducibility principles}
The governed pipeline follows five principles: deterministic execution under pinned inputs, a single geometry authority for \texttt{planView} and \texttt{elevationProfile}, artifact-first logging, offline-first validation, and causal separation between synthesis and evaluation. In practice, this means that structural evidence must remain interpretable even when CARLA is unavailable or unstable, and that later perception-facing steps are not allowed to retroactively alter the structural analysis basis.

\subsection{Artifact contract and determinism reporting}
\label{sec:artifact_contract_pipeline}
Each run preserves a compact artifact contract rather than only a final XODR file. At minimum, the governed run directory contains the input snapshot, stage outputs, key diagnostic JSON files, and run-defining hashes. This contract serves two purposes. First, it allows later stages to be audited without re-running the full conversion chain. Second, it distinguishes byte-level file drift from changes that would matter scientifically, such as altered road counts, junction counts, or reference-line geometry.

Determinism is therefore reported at the artifact level, not inferred from a single successful load. Repeated runs are compared through settings snapshots, hashes, and semantic signatures so that structural invariance can be defended even when XML serialization varies across runs \cite{peng2011reproducible,sandve2013reproducible}.

\subsection{Summary}
At Bachelor-thesis scope, the important point is not every implementation detail of the governed code base, but the staged logic of the conversion path: OSM input is sanitized, topology is repaired, semantics are enriched, geometry is stabilized, lanes and lane links are repaired, and the result is accepted only if explicit integrity gates are satisfied. This staged view is the basis for the structural experiments in Chapters~5--7 and for the simulator-dependent QA discussed in Chapter~6.
